\documentclass[11pt]{article}

\usepackage[margin=1in]{geometry}
\usepackage{amsmath,amssymb}
\usepackage{listings}
\usepackage{xcolor}
\usepackage[hidelinks]{hyperref}
\usepackage{array}

\lstset{
  basicstyle=\ttfamily\small,
  keywordstyle=\color{black}\bfseries,
  commentstyle=\color{gray},
  breaklines=true,
  frame=single,
  framerule=0.4pt,
  xleftmargin=1em,
  framexleftmargin=0.5em,
  columns=fullflexible,
  keepspaces=true,
  aboveskip=1em,
  belowskip=1em,
}

\definecolor{accentblue}{RGB}{30, 64, 175}

\lstset{
    basicstyle=\ttfamily\small,
    escapeinside={(*@}{@*)},
    frame=single,
    breaklines=true,
}

\newcommand{\added}[1]{\textcolor{accentblue}{#1}}

\title{Billing and Insurance Extensions to the FEN Schema:\\
Patient-Facing Improvements for Bills, Denials, and Coverage}
\author{Phoros Engineering}
\date{July 2026}

\begin{document}
\maketitle

\section{Overview}

This document specifies extensions to the billing and insurance facts of the
FEN schema. The current schema captures a workable claims spine: coverage
enrollment, claim submission, adjudication, payment, pre-authorization
requests and decisions, and appeals. These extensions address gaps that
prevent the schema from supporting patient-facing features such as bill
reconciliation, denial triage, deductible tracking, and structured appeals.

The changes are motivated by patient value rather than by any particular
go-to-market strategy. A patient who wants to detect a duplicate charge on a
provider bill, understand why a prior authorization was denied, see how
close they are to their deductible, or track an appeal through its
lifecycle cannot do so with the current schema. The extensions here close
those gaps.

Section~\ref{sec:institution} introduces an \texttt{Institution} entity
to replace bare strings for payers, providers, pharmacies, and other
organizations. Section~\ref{sec:provider-statement} introduces a
\texttt{ProviderStatement} payload distinct from \texttt{Claim} and
\texttt{Adjudication}. Section~\ref{sec:claim-lines} makes claim service
lines first-class. Section~\ref{sec:accumulator} adds
\texttt{PlanAccumulator} as an event-driven observation series.
Section~\ref{sec:coverage-timeline} converts \texttt{Coverage} from a
snapshot to a timeline and introduces a \texttt{PlanStructure} type
distinct from \texttt{CoverageType}. Section~\ref{sec:denial-codes}
replaces the free-text denial reason with structured CARC/RARC codes.
Section~\ref{sec:appeal-lifecycle} adds an \texttt{AppealEvent} pattern
and discusses the event-driven vs.\ state-machine design tradeoff.
Section~\ref{sec:pharmacy} adds a \texttt{PharmacyClaim} payload and a
structured \texttt{PharmacyCoverageResult}. It also represents a
point-of-prescribing formulary check as its own Fact.
Section~\ref{sec:payment-sources} expands \texttt{PayerType} and revises
\texttt{PaymentPlan} to distinguish provider-managed plans from
third-party medical financing, collections agencies, and charity
programs. Section~\ref{sec:fact-relation} adds a \texttt{FactRelation}
entity for typed connections between billing facts.

All extensions preserve the graph-oriented representation of the FEN
schema. Entities remain nodes, and cross-fact relations are represented
through typed relations and \texttt{FactId} references. The extensions use
the core \texttt{Fact} distinction between assertion provenance and evidence:
\texttt{Fact.assertion} records who asserted a billing Fact, when, and how;
\texttt{Fact.evidence} records the source citations and prior Facts that
support it.

\subsection{From a billing document to FEN Facts}

Consider a payer Explanation of Benefits (EOB) received as a PDF. The PDF is
not itself a Fact. It is preserved outside FEN as an immutable
\texttt{SourceArtifact}, whose bytes are held by a \texttt{StoredObject}.
A citation to the relevant page or region becomes a
\texttt{SourceCitationRef}. Normalization may then produce Claim and
Adjudication Facts that cite that location as evidence. The payer's claim
number is an \texttt{ExternalRef}: it identifies the source record, but does
not by itself support the normalized assertion. A typed \texttt{FactRelation}
may connect the Claim to the Adjudication. Processing state and patient-facing
reconciliation views remain \texttt{FactOperational} data or projections.

The same boundary applies to provider statements, remittance files, payer API
responses, pharmacy transactions, and Epic EHI table rows:

\begin{center}
\small
\renewcommand{\arraystretch}{1.15}
\begin{tabular}{>{\raggedright\arraybackslash}p{5.4cm} >{\raggedright\arraybackslash}p{7.6cm}}
\textbf{Concern} & \textbf{Representation} \\
\hline
Original bytes or archive & \texttt{SourceArtifact} $\rightarrow$ \texttt{StoredObject} \\
Exact supporting passage or row & \texttt{SourceCitationRef} \\
Who asserted the normalized claim & \texttt{AssertionProvenance} \\
Supporting source material or prior Facts & \texttt{EvidenceRef} \\
Source-system claim or statement identity & \texttt{ExternalRef} \\
Relationship between billing Facts & typed \texttt{FactRelation} \\
Transient processing or display state & \texttt{FactOperational} or projection \\
\end{tabular}
\end{center}

\section{Institution as a first-class entity}
\label{sec:institution}

The current schema represents payers, providers, pharmacies, and other
organizations as bare \texttt{String} fields:

\begin{lstlisting}
Coverage    { payer: String, plan: String, ... }
Claim       { payer: String, ... }
PreAuthRequest { payer: String, ... }
Encounter   { facility: Option<String>, ... }
\end{lstlisting}

Bare strings have three problems. They admit no identity resolution:
``AHN,'' ``Allegheny Health Network,'' and ``Allegheny Health Network,
Inc.''\ are three distinct string values but one institution. They
carry no canonical identifier such as an NPI (National Provider
Identifier), a Payer ID, or an NCPDP (pharmacy) number, so external
correlation is impossible. They prevent cross-fact reasoning: forty-seven
claims from the same payer appear as forty-seven distinct string values
rather than forty-seven references to one entity.

An \texttt{Institution} entity resolves this:

\begin{lstlisting}
(*@\added{Institution \{}@*)
    (*@\added{id:              InstitutionId,}@*)
    (*@\added{display\_name:    String,}@*)
    (*@\added{aliases:         Vec<String>,}@*)
    (*@\added{canonical\_ids:   Vec<CanonicalId>,}@*)
    (*@\added{institution\_type: InstitutionType,}@*)
(*@\added{\}}@*)

(*@\added{CanonicalId \{}@*)
    (*@\added{Npi(String),                   // National Provider Identifier}@*)
    (*@\added{PayerId(String),               // CMS payer identifier}@*)
    (*@\added{Ncpdp(String),                 // pharmacy}@*)
    (*@\added{TaxId(String),                 // EIN}@*)
    (*@\added{ClearinghouseId \{ system: String, value: String \},}@*)
    (*@\added{Other \{ system: String, value: String \},}@*)
(*@\added{\}}@*)

(*@\added{InstitutionType \{}@*)
    (*@\added{Payer,}@*)
    (*@\added{Provider,}@*)
    (*@\added{Pharmacy,}@*)
    (*@\added{Laboratory,}@*)
    (*@\added{ImagingCenter,}@*)
    (*@\added{Employer,}@*)
    (*@\added{ThirdPartyAdministrator,}@*)
    (*@\added{Other,}@*)
(*@\added{\}}@*)
\end{lstlisting}

\texttt{Institution} is contextual identity data, not a fourth FEN
epistemic layer. It is neither a Fact, a ProblemEpisode, nor a Narrative.
Billing Facts refer to a stable \texttt{InstitutionId}; the registry provides
canonical identity and display information. A registry merge or alias decision
must be authored and auditable, but it must not rewrite the payer or provider
name preserved in the original source artifact.

Facts reference institutions by \texttt{InstitutionId}, and the
institution registry holds display names, aliases, and canonical
identifiers. This is analogous to \texttt{SubjectId}: a stable internal
identifier that resolves to external identity systems.

The identity resolution problem is real but bounded: identity
resolution runs at ingestion, not at query time, and produces an
authored decision. When two ingestion sources refer to the same
institution differently, the resolver produces a mapping and records
the decision as an \texttt{InstitutionAliasRelation}. Ambiguous cases
surface for review rather than silently merging or duplicating.

Fields that previously carried bare institution strings are migrated to
\texttt{InstitutionId} references:

\begin{lstlisting}
Coverage {
    (*@\added{payer:      InstitutionId,}@*)
    plan:       String,
    // ...
}

Claim {
    (*@\added{payer:      InstitutionId,}@*)
    // ...
}

PreAuthRequest {
    (*@\added{payer:      InstitutionId,}@*)
    // ...
}
\end{lstlisting}

Clinical facts such as \texttt{Encounter.facility} may continue to use
bare strings in the near term, since the reasoning value of institution
identity is highest for insurance-related facts. Migration of clinical
facts to \texttt{InstitutionId} is a follow-up, not a prerequisite.

\section{ProviderStatement}
\label{sec:provider-statement}

A \texttt{ProviderStatement} represents the bill a provider sends to the
patient. It is distinct from a \texttt{Claim} (what the provider billed the
payer) and from an \texttt{Adjudication} (what the payer decided). The
provider bill and the payer's Explanation of Benefits arrive on different
timelines, use different formats, and routinely disagree. A patient
detects billing errors by comparing the two. The schema must represent
both.

\begin{lstlisting}
FactPayload {
    // ... existing variants ...

    (*@\added{ProviderStatement \{}@*)
        (*@\added{statement\_id:      String,}@*)
        (*@\added{provider:          InstitutionId,}@*)
        (*@\added{statement\_date:    Date,}@*)
        (*@\added{due\_date:          Option<Date>,}@*)
        (*@\added{service\_period:    TemporalAnchor,}@*)
        (*@\added{lines:             Vec<StatementLine>,}@*)
        (*@\added{total\_charges:     Money,}@*)
        (*@\added{insurance\_paid:    Money,}@*)
        (*@\added{adjustments:       Money,}@*)
        (*@\added{prior\_balance:     Money,}@*)
        (*@\added{amount\_due:        Money,}@*)
        (*@\added{related\_claim:     Option<FactId>,}@*)
    (*@\added{\},}@*)
}
\end{lstlisting}

The \texttt{related\_claim} field links to the \texttt{Claim} fact this
statement corresponds to, when the correspondence can be established. A
missing link is itself informative: it means the patient has received a
bill for a service that has no matching claim in their record, which
often indicates the provider has not yet billed insurance.

\begin{lstlisting}
(*@\added{StatementLine \{}@*)
    (*@\added{service\_date:         Date,}@*)
    (*@\added{description:          String,}@*)
    (*@\added{procedure\_code:       Option<CodedValue>, // CPT or HCPCS}@*)
    (*@\added{units:                u32,}@*)
    (*@\added{charge:               Money,}@*)
    (*@\added{insurance\_payment:    Money,}@*)
    (*@\added{adjustment:           Money,}@*)
    (*@\added{patient\_responsibility: Money,}@*)
(*@\added{\}}@*)
\end{lstlisting}

An EOB received as a PDF or 835 transaction is preserved as a
\texttt{SourceArtifact}; it is not recoverable merely by re-rendering the
normalized Adjudication Facts. A later \texttt{ExplanationOfBenefits} Fact
payload would be justified only if it captured an EOB-level domain assertion,
not simply because the original file needs a container. Either way, the
original artifact and the citations used during normalization remain intact.

\section{Claim service lines as first-class structure}
\label{sec:claim-lines}

The current \texttt{Claim} payload treats billed codes and billed amounts
as aggregate:

\begin{lstlisting}
Claim {
    claim_id:       String,
    billed_codes:   Vec<CodedValue>,
    billed_amount:  Money,
    payer:          String,
}
\end{lstlisting}

Real claims have multiple service lines, each with its own procedure
code, unit count, billed amount, and adjudication outcome. To detect
duplicate charges (the same code billed twice), unbundling violations
(services billed separately that should have been billed together), or
upcoding (a higher-severity code than the encounter warranted), a claim
must expose its lines. Similarly, denials frequently occur at the line
level: a claim as a whole may be marked \texttt{PartiallyPaid} while one
line is denied and the others are paid.

\begin{lstlisting}
Claim {
    claim_id:          String,
    (*@\added{lines:             Vec<ClaimLine>,}@*)
    (*@\added{total\_billed:      Money, // sum of line.billed\_amount}@*)
    (*@\added{payer:             InstitutionId,}@*)
    (*@\added{claim\_type:        ClaimType, // Professional | Institutional}@*)
    (*@\added{service\_period:    TemporalAnchor,}@*)
}
\end{lstlisting}

The aggregate \texttt{billed\_codes} field is removed; consumers derive
the code list from \texttt{lines}. The aggregate \texttt{billed\_amount}
is renamed \texttt{total\_billed} and constrained to equal the sum of
line-level billed amounts. Ingestion validates this invariant.

\begin{lstlisting}
(*@\added{ClaimLine \{}@*)
    (*@\added{line\_number:     u32,}@*)
    (*@\added{service\_date:    Date,}@*)
    (*@\added{procedure\_code:  CodedValue,        // CPT or HCPCS}@*)
    (*@\added{modifiers:       Vec<CodedValue>,   // CPT modifiers}@*)
    (*@\added{diagnosis\_refs:  Vec<CodedValue>,   // ICD-10 pointers}@*)
    (*@\added{units:           u32,}@*)
    (*@\added{billed\_amount:   Money,}@*)
    (*@\added{place\_of\_service: Option<CodedValue>,}@*)
    (*@\added{rendering\_provider: Option<String>,}@*)
(*@\added{\}}@*)
\end{lstlisting}

\texttt{Adjudication} correspondingly gains line-level detail:

\begin{lstlisting}
Adjudication {
    claim_id:                 String,
    (*@\added{lines:                    Vec<AdjudicationLine>,}@*)
    (*@\added{total\_allowed:            Money,}@*)
    (*@\added{total\_patient\_resp:       Money,}@*)
    status:                   AdjudicationStatus,
    (*@\added{processed\_date:           Date,}@*)
}

(*@\added{AdjudicationLine \{}@*)
    (*@\added{line\_number:          u32, // matches ClaimLine.line\_number}@*)
    (*@\added{allowed\_amount:       Money,}@*)
    (*@\added{paid\_amount:          Money,}@*)
    (*@\added{patient\_responsibility: Money,}@*)
    (*@\added{deductible\_applied:   Money,}@*)
    (*@\added{coinsurance\_applied:  Money,}@*)
    (*@\added{copay\_applied:        Money,}@*)
    (*@\added{denial\_codes:         Vec<CodedValue>, // CARC/RARC per Sec.~\ref{sec:denial-codes}}@*)
    (*@\added{line\_status:          AdjudicationStatus,}@*)
(*@\added{\}}@*)
\end{lstlisting}

\section{PlanAccumulator}
\label{sec:accumulator}

Deductible and out-of-pocket accumulator state is derived from
adjudications but reported by the payer as a snapshot. Patients need to
see it directly. Providers use it to quote patient responsibility before
a service. The current schema has no representation of accumulator state.

\begin{lstlisting}
FactPayload {
    // ...

    (*@\added{PlanAccumulator \{}@*)
        (*@\added{coverage\_ref:              FactId,   // link to Coverage fact}@*)
        (*@\added{as\_of:                     Date,}@*)
        (*@\added{plan\_year:                 u16,}@*)
        (*@\added{individual\_deductible:     AccumulatorAmount,}@*)
        (*@\added{family\_deductible:         Option<AccumulatorAmount>,}@*)
        (*@\added{individual\_oop\_max:        AccumulatorAmount,}@*)
        (*@\added{family\_oop\_max:            Option<AccumulatorAmount>,}@*)
        (*@\added{network\_scope:             NetworkScope,}@*)
        (*@\added{reported\_via:              Option<AccumulatorChannel>,}@*)
    (*@\added{\},}@*)
}

(*@\added{AccumulatorAmount \{}@*)
    (*@\added{limit:  Money,}@*)
    (*@\added{met:    Money,       // amount accumulated toward the limit}@*)
(*@\added{\}}@*)

(*@\added{NetworkScope \{ InNetwork, OutOfNetwork, Combined \}}@*)

(*@\added{AccumulatorChannel \{}@*)
    (*@\added{PayerApi,          // 270/271 eligibility or Patient Access API}@*)
    (*@\added{ProviderPortal,    // captured at time of service}@*)
    (*@\added{PatientStatement,  // from paper or portal EOB}@*)
(*@\added{\}}@*)
\end{lstlisting}

\texttt{reported\_via} describes a domain communication channel; it is not
evidence provenance. The enclosing Fact cites the response, statement, or
table row through \texttt{Fact.evidence}. A value computed from prior
Adjudication Facts instead uses \texttt{AssertionMethod::Inferred}, cites
those Facts through \texttt{EvidenceRef::PriorFact}, and records the evaluator
or rule version used.

\textbf{Design decision: event-driven, not superseded.} Each
\texttt{PlanAccumulator} fact is an independent observation, not a
correction of a prior one. A new accumulator reading does not supersede
the previous one; both facts remain \texttt{Active} and reflect the
accumulator state at their respective \texttt{as\_of} dates. This
diverges from an earlier design that used
\texttt{FactStatus::Superseded} with an \texttt{AccumulatorUpdate}
supersession reason.

The rationale: supersession is appropriate when a fact is being
corrected or refined --- the earlier fact was wrong or partial, and the
newer fact replaces it. Accumulator readings are not corrections. Each
reading is a true observation at a real point in time, and the
progression of readings is itself informative. A patient who hit their
deductible at three different visits wants to see the progression
(``\$500 met in March, \$1{,}200 met in May, \$3{,}000 met in July''),
not just the most recent snapshot. Supersession would erase the
progression.

This mirrors how \texttt{LabResult} works in the existing schema. Every
lab draw is a separate \texttt{LabResult} fact with its own timestamp,
not a supersession of the previous draw. A hemoglobin value from last
month is not made wrong by a hemoglobin value from this month; both
are true observations at their respective times. The same reasoning
applies to accumulators.

The read model selects the most recent \texttt{PlanAccumulator} for the
relevant coverage and plan year. A historical query such as ``what was my
deductible progress on July 1?'' selects the latest observation available at
that date. Neither query requires supersession.

Supersession is still available for the rare case where a payer
retroactively corrects an accumulator reading. The prior Fact is then marked
\texttt{Superseded} for clinical refinement. This treats the change as a
correction of one specific observation rather than as a routine update.

\section{Coverage timeline}
\label{sec:coverage-timeline}

The current \texttt{Coverage} payload represents coverage as a fact
without explicit temporal bounds:

\begin{lstlisting}
Coverage {
    payer:          String,
    plan:           String,
    member_id:      String,
    group_id:       Option<String>,
    coverage_type:  CoverageType,
}
\end{lstlisting}

Coverage changes: patients change employers, plans renew, dependents are
added, secondary coverage begins. Coverage must be represented as a
period, not a point. The \texttt{Fact.occurred\_at} field can carry a
\texttt{TemporalAnchor::Period}, but the schema does not currently
express the domain-specific fields that make coverage-as-timeline
usable.

\begin{lstlisting}
Coverage {
    (*@\added{payer:              InstitutionId,}@*)
    plan:               String,
    member_id:          String,
    group_id:           Option<String>,
    coverage_type:      CoverageType,
    (*@\added{plan\_structure:     PlanStructure,}@*)
    (*@\added{effective\_from:     Date,}@*)
    (*@\added{effective\_through:  Option<Date>,   // None means currently active}@*)
    (*@\added{plan\_year:          Option<u16>,}@*)
    (*@\added{network\_status:     NetworkStatus,}@*)
    (*@\added{relationship\_to\_subscriber: SubscriberRelationship,}@*)
    (*@\added{policy\_holder:      Option<String>,}@*)
}

(*@\added{NetworkStatus \{ InNetwork, OutOfNetwork, Unknown \}}@*)

(*@\added{SubscriberRelationship \{}@*)
    (*@\added{Self\_,     // patient is the subscriber}@*)
    (*@\added{Spouse,}@*)
    (*@\added{Dependent,}@*)
    (*@\added{Other,}@*)
(*@\added{\}}@*)
\end{lstlisting}

\texttt{effective\_through: None} indicates the coverage is currently
active. When coverage ends, the fact is superseded and a new
\texttt{Coverage} fact is created for the replacement coverage.

The \texttt{plan\_year} field distinguishes ``same employer, different
plan year'' (a plan renewal with updated benefits) from ``different
employer.'' This distinction matters for accumulator resets: at the
start of a new plan year the deductible resets even if the coverage is
continuous.

\subsection{PlanStructure}
\label{sec:plan-structure}

The existing \texttt{CoverageType} enum
(\texttt{Commercial\ \textbar{}\ Medicare\ \textbar{}\ Medicaid\ \textbar{}\ Tricare\ \textbar{}\ SelfPay\ \textbar{}\ Other})
distinguishes \emph{who pays} but not \emph{how the plan is
structured}. A commercial HDHP with a \$6{,}000 deductible and 20\%
coinsurance behaves totally differently from a commercial PPO with
\$30 copays and no deductible, and the patient's cost-share reasoning
depends on the structure. \texttt{PlanStructure} captures the
structural distinction:

\begin{lstlisting}
(*@\added{PlanStructure \{}@*)
    (*@\added{plan\_type:               PlanType,}@*)
    (*@\added{has\_deductible:          bool,}@*)
    (*@\added{has\_coinsurance:         bool,}@*)
    (*@\added{has\_copay:               bool,}@*)
    (*@\added{is\_hsa\_eligible:         bool,   // HDHP paired with HSA}@*)
(*@\added{\}}@*)

(*@\added{PlanType \{}@*)
    (*@\added{Hmo,}@*)
    (*@\added{Ppo,}@*)
    (*@\added{Epo,}@*)
    (*@\added{Pos,}@*)
    (*@\added{Hdhp,      // High-Deductible Health Plan}@*)
    (*@\added{Indemnity,}@*)
    (*@\added{Unknown,}@*)
(*@\added{\}}@*)
\end{lstlisting}

\texttt{PlanStructure} is deliberately minimal. It does not attempt to
model the full plan design (formulary, service-level coverage rules,
network tiers, benefit-year rules). Those belong in a separate
plan-design contextual model outside \texttt{FactPayload}. The
structural type is included on \texttt{Coverage} because the patient
needs to know whether their next visit will hit a deductible or a
copay, and that requires knowing the plan structure.

\texttt{PlanStructure} is a field on \texttt{Coverage} rather than a
separate fact. It describes the coverage and does not vary
independently over time --- when the plan structure changes (e.g., the
employer switches plans mid-year), a new \texttt{Coverage} fact is
created and the prior one is superseded.

\section{Structured denial reasons}
\label{sec:denial-codes}

The current \texttt{PreAuthDecision} payload uses a free-text denial
reason:

\begin{lstlisting}
PreAuthDecision {
    authorization_id:   String,
    status:             AuthorizationStatus,
    reason:             Option<String>,
    approved_service:   Option<CodedValue>,
    valid_through:      Option<Date>,
}
\end{lstlisting}

Free-text denial reasons are unparseable. Payers actually communicate
denials using standardized codes: CARC (Claim Adjustment Reason Codes,
maintained by X12) and RARC (Remittance Advice Remark Codes, maintained
by CMS). Structured codes enable a real feature: given a denial reason,
determine whether the patient has a viable appeal path, what
documentation is required, and what the timeline is.

\begin{lstlisting}
PreAuthDecision {
    authorization_id:   String,
    status:             AuthorizationStatus,
    (*@\added{denial\_codes:       Vec<CodedValue>,   // CARC and RARC}@*)
    (*@\added{denial\_narrative:   Option<String>,    // free-text supplement}@*)
    approved_service:   Option<CodedValue>,
    valid_through:      Option<Date>,
    (*@\added{decision\_date:      Date,}@*)
    (*@\added{appeal\_deadline:    Option<Date>,}@*)
}
\end{lstlisting}

The same treatment applies to \texttt{AdjudicationLine.denial\_codes}
(see Section~\ref{sec:claim-lines}). CARC and RARC vocabularies are
represented as \texttt{CodedValue} with the code system identifier
distinguishing the two.

The free-text \texttt{denial\_narrative} field is preserved as an
optional supplement. Payer notices frequently include narrative
explanations that add context beyond the structured code (``member has
not tried step-therapy tier 1 medication'' is a RARC N123 with a
specific meaning; the narrative provides the humanized version).

When a payer supplies the codes, the normalized Fact cites the relevant notice
or transaction location. When Phoros maps free text to codes, that mapping is
a separate inferred coding assertion supported by the original denial Fact.
It does not silently convert model output into payer-authored content.

\section{Appeal lifecycle}
\label{sec:appeal-lifecycle}

The current \texttt{Appeal} payload captures the case but not its
progression:

\begin{lstlisting}
Appeal {
    original_decision_id:   FactId,
    grounds:                String,
    supporting_facts:       Vec<FactId>,
}
\end{lstlisting}

An appeal is a multi-week process with distinct states: filed,
acknowledged, additional information requested, decided at first level,
escalated to second level, escalated to external review, resolved.
Without lifecycle tracking, the schema cannot answer basic questions
such as ``is this appeal still open,'' ``did it succeed,'' or ``what
level is it at.''

\subsection{Design decision: event-driven vs.\ state-machine}

Two designs were considered for appeal lifecycle. The adopted design uses the
same event pattern as medication modeling: an \texttt{Appeal} Fact establishes
the case, and \texttt{AppealEvent} Facts record each state transition. The
current state of an appeal is computed by folding over its events. The
state-machine design (not adopted) would place a mutable \texttt{status}
field on \texttt{Appeal} that gets updated as the appeal progresses,
without separate event facts.

\paragraph{Arguments for event-driven (adopted).}
Consistent with the medication-event idiom, so
readers of the schema encounter one pattern rather than two. Each
\texttt{AppealEvent} preserves its own assertion metadata and evidence: who
reported the event, when and how it was asserted, and which document or prior
Fact supports it. The full
history is preserved as an immutable sequence of facts rather than
compressed into a single mutable status. Payer decisions, patient
escalations, and external review outcomes are genuine events with
external actors, dates, and documents, and treating them as first-class
facts respects that. The read model for ``current state'' is trivially
implementable and cacheable.

\paragraph{Arguments for state-machine (not adopted).}
Simpler for common queries such as ``list all open appeals'' --- a
single field lookup rather than a fold over events. Fewer total facts.
The FEN schema already has a state-machine mechanism in
\texttt{FactStatus} (\texttt{Active}, \texttt{Superseded},
\texttt{EnteredInError}), so appeal state could be modeled the same way
rather than introducing a second pattern.

\paragraph{Why event-driven wins here.}
An \texttt{AppealEvent} corresponds to a specific real-world
communication: a letter from the payer acknowledging the appeal, a
request for additional documentation, a decision letter with a
determination, a notice of escalation to external review. The appeal
itself is a case; the events are what came in the mail. Treating those
communications as first-class facts (with dates, senders, assertion metadata,
and evidence) is faithful to how the appeal actually unfolds. A
mutable \texttt{status} field cannot carry that structure without
either losing information or duplicating it in a shadow event log.

The state-machine design is preferable when the entity being tracked
is genuinely a single object whose state changes internally, and the
history of state changes is not itself an object of interest. Neither
condition holds for appeals.

The same reasoning was applied to \texttt{PlanAccumulator}
(Section~\ref{sec:accumulator}), where each observation is a first-class
event and supersession was rejected.

\subsection{Schema}

The pattern mirrors medication-order event modeling. The \texttt{Appeal} Fact
establishes the case; \texttt{AppealEvent} Facts track state transitions.

\begin{lstlisting}
Appeal {
    (*@\added{appeal\_id:              Option<String>, // source-system ID}@*)
    original_decision_id:   FactId,
    (*@\added{appeal\_level:           AppealLevel,}@*)
    grounds:                String,
    (*@\added{filed\_date:             Date,}@*)
    (*@\added{filed\_by:               AppealFiler,}@*)
}

(*@\added{AppealLevel \{}@*)
    (*@\added{Internal1,        // first-level internal payer review}@*)
    (*@\added{Internal2,        // second-level internal review}@*)
    (*@\added{ExternalReview,   // independent external review}@*)
    (*@\added{StateComplaint,   // filed with state insurance commissioner}@*)
(*@\added{\}}@*)

(*@\added{AppealFiler \{ Patient, Provider, AuthorizedRepresentative \}}@*)

(*@\added{AppealEvent \{}@*)
    (*@\added{appeal\_ref:      FactId,   // links to the Appeal fact}@*)
    (*@\added{event\_type:      AppealEventType,}@*)
    (*@\added{notes:           Option<String>,}@*)
(*@\added{\},}@*)

(*@\added{AppealEventType \{}@*)
    (*@\added{Acknowledged \{ acknowledged\_date: Date \},}@*)
    (*@\added{InfoRequested \{ items: Vec<String>, deadline: Date \},}@*)
    (*@\added{InfoSubmitted \{ items: Vec<FactId> \},}@*)
    (*@\added{Decided \{}@*)
        (*@\added{outcome:         AppealOutcome,}@*)
        (*@\added{decision\_date:   Date,}@*)
        (*@\added{denial\_codes:    Vec<CodedValue>, // if upheld}@*)
    (*@\added{\},}@*)
    (*@\added{Escalated \{ to\_level: AppealLevel \},}@*)
    (*@\added{Withdrawn \{ withdrawn\_date: Date, reason: Option<String> \},}@*)
(*@\added{\}}@*)

(*@\added{AppealOutcome \{ Overturned, PartiallyOverturned, Upheld \}}@*)
\end{lstlisting}

The enclosing Appeal and AppealEvent Facts use \texttt{Fact.evidence} for
supporting documents and prior Facts. A document is cited with
\texttt{EvidenceRef::Source}; a prior clinical or billing Fact is cited with
\texttt{EvidenceRef::PriorFact}. The \texttt{InfoSubmitted.items} field
remains because it describes the content of that event---which items were
submitted---rather than the epistemic support for the event record itself.

Both \texttt{Appeal} and \texttt{AppealEvent} become \texttt{FactPayload}
variants. Corresponding \texttt{FactType} variants are added:
\texttt{FactType::AppealEvent}.

\section{PharmacyClaim}
\label{sec:pharmacy}

The current \texttt{Claim} payload is designed for medical claims
(CPT/HCPCS codes, professional or institutional claim types). Pharmacy
claims use different vocabularies and pricing structures: NDC codes for
drug identification, days-supply and quantity dispensed, DAW (Dispense
As Written) codes, formulary tier, and PBM adjudication that includes
rebates and MAC (Maximum Allowable Cost) pricing. Forcing pharmacy
claims into the medical \texttt{Claim} structure loses information and
prevents useful features (formulary alternative suggestions, generic
substitution tracking, PBM cost comparison).

A separate \texttt{PharmacyClaim} payload:

\begin{lstlisting}
FactPayload {
    // ...

    (*@\added{PharmacyClaim \{}@*)
        (*@\added{rx\_claim\_id:         String,}@*)
        (*@\added{medication\_order\_ref: Option<FactId>, // link to medication-order Fact}@*)
        (*@\added{ndc:                 CodedValue,   // National Drug Code}@*)
        (*@\added{drug\_name:           String,}@*)
        (*@\added{days\_supply:         u32,}@*)
        (*@\added{quantity\_dispensed:  Decimal,}@*)
        (*@\added{unit:                CodedValue,}@*)
        (*@\added{dispense\_date:       Date,}@*)
        (*@\added{pharmacy:            InstitutionId,}@*)
        (*@\added{prescriber:          String,}@*)
        (*@\added{daw\_code:            Option<u8>,}@*)
        (*@\added{payer:               InstitutionId,}@*)
        (*@\added{pbm:                 Option<InstitutionId>,}@*)
        (*@\added{billed\_amount:       Money,}@*)
        (*@\added{plan\_paid:           Money,}@*)
        (*@\added{patient\_paid:        Money,}@*)
        (*@\added{coverage\_result:     PharmacyCoverageResult,}@*)
    (*@\added{\},}@*)
}
\end{lstlisting}

The \texttt{medication\_order\_ref} field links the pharmacy claim back to
the originating medication-order Fact when available. This link is
essential for reasoning about medication adherence (comparing prescribed
days-supply to actual dispensing intervals) and for detecting formulary
switches (the pharmacy dispensed a different NDC than prescribed).

\subsection{PharmacyCoverageResult}
\label{sec:pharmacy-coverage-result}

An earlier design placed covered, not-covered, prior-authorization,
step-therapy, quantity-limit, and cash-pay outcomes in one status enum. In
practice
pharmacy rejections often carry multiple reasons simultaneously
(``not covered on formulary, and also requires step therapy, and also
has a quantity limit''), which an enum with mutually exclusive variants
cannot represent. A struct that composes the individual conditions
handles this better:

\begin{lstlisting}
(*@\added{PharmacyCoverageResult \{}@*)
    (*@\added{outcome:                CoverageOutcome,}@*)
    (*@\added{tier:                   Option<u8>,}@*)
    (*@\added{prior\_auth\_required:    bool,}@*)
    (*@\added{step\_therapy\_required:  bool,}@*)
    (*@\added{quantity\_limit\_hit:     bool,}@*)
    (*@\added{suggested\_alternative:  Option<CodedValue>, // NDC of covered alt}@*)
    (*@\added{rejection\_codes:        Vec<CodedValue>,    // NCPDP reject codes}@*)
    (*@\added{requires\_manufacturer\_assistance: bool,}@*)
(*@\added{\}}@*)

(*@\added{CoverageOutcome \{}@*)
    (*@\added{Paid,}@*)
    (*@\added{Rejected,}@*)
    (*@\added{Pending,}@*)
    (*@\added{CashPay,     // adjudicated as cash by patient choice or plan denial}@*)
(*@\added{\}}@*)
\end{lstlisting}

\subsection{MedicationCoverageCheck at prescribing}
\label{sec:formulary-tier}

Formulary tier at the point of prescribing is separately useful from
formulary tier at adjudication. Provider EHRs increasingly run
real-time benefit checks (RTBC) via NCPDP formulary APIs at the moment
a prescription is written, so the tier and prior-auth requirements
are known before the patient ever reaches the pharmacy.

The RTBC result must not be inserted into the clinician-authored medication
order as though the clinician asserted it. The PBM or benefit service made a
separate assertion, often from a separate response artifact. The adopted
representation is therefore another Fact:

\begin{lstlisting}
FactPayload {
    (*@\added{MedicationCoverageCheck \{}@*)
        (*@\added{medication:          CodedValue,}@*)
        (*@\added{medication\_order\_ref: Option<FactId>,}@*)
        (*@\added{coverage\_ref:       Option<FactId>,}@*)
        (*@\added{checked\_at:         Timestamp,}@*)
        (*@\added{formulary\_tier:     Option<u8>,}@*)
        (*@\added{formulary\_status:   FormularyStatus,}@*)
    (*@\added{\},}@*)
}

(*@\added{FormularyStatus \{}@*)
    (*@\added{Covered,}@*)
    (*@\added{NotCovered,}@*)
    (*@\added{PriorAuthRequired,}@*)
    (*@\added{StepTherapyRequired,}@*)
    (*@\added{QuantityLimited,}@*)
    (*@\added{NotChecked,     // no RTBC was performed}@*)
(*@\added{\}}@*)
\end{lstlisting}

The enclosing Fact records the benefit service's assertion metadata and cites
the RTBC response through \texttt{Fact.evidence}. Its
\texttt{medication\_order\_ref} is a semantic connection to the order, not a
substitute for that evidence. The \texttt{MedicationCoverageCheck} represents
``what the real-time benefit check said at prescribing.'' The
\texttt{PharmacyCoverageResult} on \texttt{PharmacyClaim} represents
``what the PBM actually adjudicated at dispensing.'' The two often
agree; when they diverge, the divergence itself is informative
(``the RTBC said covered, but the PBM rejected --- something changed
between prescribing and pickup'').

Formulary is properly a plan-design property, and the tier value only
has meaning in the context of a specific plan's formulary. Modeling
the plan formulary itself is out of scope for this document
(see Section~\ref{sec:what-not-added}); capturing the tier at
observation time is sufficient for patient-facing reasoning.

\section{Payment sources and methods}
\label{sec:payment-sources}

The current \texttt{Payment} payload and \texttt{PayerType} enum treat
payments as a small closed set:

\begin{lstlisting}
Payment {
    claim_id:  String,
    amount:    Money,
    paid_by:   PayerType,
}

PayerType {
    PrimaryInsurance,
    SecondaryInsurance,
    Patient,
    Other,
}
\end{lstlisting}

Real patient payments come from a wider set of sources with different
tax and financial-management implications. HSA/FSA payments are pre-tax
dollars with reporting requirements. Payment plans are partial payments
over time and change what the patient owes. Medical credit lines
(CareCredit) are debt with interest terms. Charity care and financial
assistance reduce the patient's obligation. Collections activity
indicates a payment failure that affects credit. Each of these matters
to the patient and to any downstream reasoning about what the patient
actually spent on healthcare.

\begin{lstlisting}
Payment {
    claim_id:              String,
    (*@\added{provider\_statement\_ref: Option<FactId>, // link to statement paid}@*)
    amount:                Money,
    paid_by:               PayerType,
    (*@\added{payment\_method:        PaymentMethod,}@*)
    (*@\added{payment\_date:          Date,}@*)
    (*@\added{payment\_plan\_ref:      Option<FactId>,}@*)
}

(*@\added{PayerType \{}@*)
    (*@\added{PrimaryInsurance,}@*)
    (*@\added{SecondaryInsurance,}@*)
    (*@\added{Medicare,}@*)
    (*@\added{Medicaid,}@*)
    (*@\added{Patient,}@*)
    (*@\added{HSA,}@*)
    (*@\added{FSA,}@*)
    (*@\added{HRA,}@*)
    (*@\added{MedicalCredit \{ issuer: String \},        // e.g. CareCredit}@*)
    (*@\added{CharityCare \{ program: String \},}@*)
    (*@\added{FinancialAssistance \{ program: String \},}@*)
    (*@\added{Collections \{ agency: String \},}@*)
    (*@\added{WorkersComp,}@*)
    (*@\added{AutoInsurance,}@*)
    (*@\added{Other \{ description: String \},}@*)
(*@\added{\}}@*)

(*@\added{PaymentMethod \{}@*)
    (*@\added{Card,}@*)
    (*@\added{Check,}@*)
    (*@\added{BankTransfer,}@*)
    (*@\added{Cash,}@*)
    (*@\added{PaymentPlan,}@*)
    (*@\added{Adjustment,       // contractual or courtesy write-off}@*)
    (*@\added{Other,}@*)
(*@\added{\}}@*)
\end{lstlisting}

A payment plan itself is not a payment; it is an agreement that
produces a sequence of payments. Payment plans need their own payload.

An earlier design assumed the provider always manages the payment plan. Some
plans instead use third-party medical financing companies, including
CareCredit, PayZen, AccessOne, and ClearBalance. Others are managed by
collections agencies, hospital charity programs, or state assistance. The
\texttt{PaymentPlan} payload uses a \texttt{PaymentPlanManager} enum
to distinguish these:

\begin{lstlisting}
FactPayload {
    // ...

    (*@\added{PaymentPlan \{}@*)
        (*@\added{plan\_id:            String,}@*)
        (*@\added{manager:            PaymentPlanManager,}@*)
        (*@\added{total\_amount:       Money,}@*)
        (*@\added{monthly\_amount:     Money,}@*)
        (*@\added{start\_date:         Date,}@*)
        (*@\added{end\_date:           Option<Date>,}@*)
        (*@\added{interest\_rate:      Option<Decimal>,}@*)
        (*@\added{related\_statements: Vec<FactId>,  // statements covered by plan}@*)
    (*@\added{\},}@*)
}

(*@\added{PaymentPlanManager \{}@*)
    (*@\added{Provider \{ institution: InstitutionId \},}@*)
    (*@\added{MedicalFinancing \{ issuer: String \},   // e.g. CareCredit, PayZen}@*)
    (*@\added{CollectionsAgency \{ agency: String \},}@*)
    (*@\added{CharityProgram \{ program: String \},   // hospital or nonprofit}@*)
    (*@\added{StateAssistance \{ program: String \},}@*)
    (*@\added{Other \{ description: String \},}@*)
(*@\added{\}}@*)
\end{lstlisting}

Individual \texttt{Payment} facts made under a payment plan link back
via \texttt{payment\_plan\_ref}. This lets the patient see, on any
statement, whether the balance is being paid down under an active
agreement, whether it is with a third-party financier accruing interest,
or whether it has fallen into collections.

Earlier \texttt{PayerType} variants such as \texttt{MedicalCredit} and
\texttt{Collections} conflated two questions. \texttt{PayerType} answers
\emph{what funded this payment}; \texttt{PaymentPlanManager} answers
\emph{who administers the agreement}. For a patient paying a CareCredit
balance, the Payment Fact records the patient as the funding source, while the
related PaymentPlan records CareCredit as the financing manager. The two
dimensions are orthogonal.

\section{FactRelation for billing}
\label{sec:fact-relation}

Billing Facts currently refer to one another through inline \texttt{FactId}
fields. Examples include:

\begin{itemize}
\item \texttt{Adjudication.claim\_id} and \texttt{Payment.claim\_id};
\item \texttt{Appeal.original\_decision\_id} and \texttt{AppealEvent.appeal\_ref};
\item \texttt{Payment.provider\_statement\_ref} and \texttt{Payment.payment\_plan\_ref};
\item \texttt{PlanAccumulator.coverage\_ref} and \texttt{PharmacyClaim.medication\_order\_ref};
\item the order and coverage references on \texttt{MedicationCoverageCheck};
\item \texttt{ProviderStatement.related\_claim} and \texttt{PaymentPlan.related\_statements}.
\end{itemize}

These are implicit relations. The existing FEN schema handles
clinical relations through \texttt{EpisodeRelation} and
\texttt{NarrativeRelation} --- explicit relation types with their own
assertion metadata. For consistency, and because billing has genuinely
many-to-many relations that inline \texttt{FactId} fields cannot
represent (one payment covering multiple claims; one claim generating
multiple statements over time), a \texttt{FactRelation} entity is introduced.

\begin{lstlisting}
(*@\added{FactRelation \{}@*)
    (*@\added{id:              RelationId,}@*)
    (*@\added{from\_fact:       FactId,}@*)
    (*@\added{to\_fact:         FactId,}@*)
    (*@\added{relation\_type:   FactRelationType,}@*)
    (*@\added{assertion:       AssertionProvenance,}@*)
    (*@\added{evidence:        Vec<EvidenceRef>,}@*)
(*@\added{\}}@*)

(*@\added{FactRelationType \{}@*)
    (*@\added{// Billing relations}@*)
    (*@\added{ClaimAdjudicatedBy,         // Claim -> Adjudication}@*)
    (*@\added{ClaimPaidBy,                // Claim -> Payment}@*)
    (*@\added{StatementBillsClaim,        // ProviderStatement -> Claim}@*)
    (*@\added{StatementPaidBy,            // ProviderStatement -> Payment}@*)
    (*@\added{AppealsDecision,            // Appeal -> PreAuthDecision or Adjudication}@*)
    (*@\added{AppealEventBelongsTo,       // AppealEvent -> Appeal}@*)
    (*@\added{PaymentPlanCovers,          // PaymentPlan -> ProviderStatement}@*)
    (*@\added{PaymentUnderPlan,           // Payment -> PaymentPlan}@*)
    (*@\added{AccumulatorForCoverage,     // PlanAccumulator -> Coverage}@*)
    (*@\added{PharmacyClaimFulfillsOrder, // PharmacyClaim -> medication-order Fact}@*)
    (*@\added{CoverageCheckForOrder,      // MedicationCoverageCheck -> medication-order Fact}@*)
    (*@\added{CoverageCheckForCoverage,   // MedicationCoverageCheck -> Coverage}@*)
    (*@\added{PreAuthForService,          // PreAuthRequest -> planned service reference}@*)
    (*@\added{PreAuthDecisionOnRequest,   // PreAuthDecision -> PreAuthRequest}@*)
    // ... additional types added as needed
(*@\added{\}}@*)
\end{lstlisting}

The direction is part of each relation type's contract. For example,
\texttt{StatementBillsClaim} always points from a ProviderStatement Fact to a
Claim Fact, while \texttt{ClaimPaidBy} always points from a Claim Fact to a
Payment Fact. Ingestion validates the permitted endpoint types.

\texttt{FactRelation} is for domain structure, not ordinary evidentiary
support. If prior Facts support an Appeal Fact, the Appeal records them in
\texttt{Fact.evidence}. A separate ``supports appeal'' edge would duplicate
that canonical evidence relationship.

\paragraph{Migration of inline \texttt{FactId} fields.}
Existing inline references are preserved during a transition period.
The intent is that ingestion produces both the inline field (for
backward compatibility with existing queries) and the corresponding
\texttt{FactRelation} (for the graph representation). Consumers migrate
to querying via relations over time. Once no consumer depends on the
inline fields, they are removed in a subsequent schema revision.
Fields identified for removal:

\begin{itemize}
\item \texttt{Adjudication.claim\_id} $\rightarrow$ \texttt{ClaimAdjudicatedBy}
\item \texttt{Payment.claim\_id} $\rightarrow$ \texttt{ClaimPaidBy}
\item \texttt{Payment.provider\_statement\_ref} $\rightarrow$ \texttt{StatementPaidBy}
\item \texttt{Payment.payment\_plan\_ref} $\rightarrow$ \texttt{PaymentUnderPlan}
\item \texttt{Appeal.original\_decision\_id} $\rightarrow$ \texttt{AppealsDecision}
\item \texttt{Appeal.supporting\_facts} $\rightarrow$ \texttt{Fact.evidence}
\item \texttt{AppealEvent.appeal\_ref} $\rightarrow$ \texttt{AppealEventBelongsTo}
\item \texttt{PlanAccumulator.coverage\_ref} $\rightarrow$ \texttt{AccumulatorForCoverage}
\item \texttt{PharmacyClaim.medication\_order\_ref} $\rightarrow$ \texttt{PharmacyClaimFulfillsOrder}
\item \texttt{MedicationCoverageCheck.medication\_order\_ref} $\rightarrow$ \texttt{CoverageCheckForOrder}
\item \texttt{MedicationCoverageCheck.coverage\_ref} $\rightarrow$ \texttt{CoverageCheckForCoverage}
\item \texttt{ProviderStatement.related\_claim} $\rightarrow$ \texttt{StatementBillsClaim}
\item \texttt{PaymentPlan.related\_statements} $\rightarrow$ \texttt{PaymentPlanCovers}
\end{itemize}

\paragraph{Advantages.}
Relations carry assertion metadata and evidence, so ``the payer told us this
claim goes with this statement'' is distinguishable from ``Phoros inferred
this relation because the amounts and dates match.'' Many-to-many relations
are directly representable. The graph
becomes uniformly queryable: ``what facts relate to this appeal''
is a single relation lookup, not a scan across multiple inline fields
on multiple fact types. New relation types are added without changing
existing fact structures.

\paragraph{Costs.}
More entities to store, and joins are required to reconstruct related
facts. The database engine must efficiently query relations by
\texttt{from\_fact} and \texttt{to\_fact}, which requires appropriate
indexing. Ingestion pipelines produce more entities per source event.
These costs are the standard tradeoff of graph-oriented modeling and
are consistent with the tradeoffs already accepted for
\texttt{EpisodeRelation} and \texttt{NarrativeRelation}.

\section{Updated FactType enumeration}

The \texttt{FactType} enum expands to include the new variants:

\begin{lstlisting}
pub enum FactType {
    Measurement,
    Prescription,
    Administration,
    Procedure,
    Diagnosis,
    ImagingResult,
    LabResult,
    Note,
    Encounter,
    Referral,
    Document,
    Coverage,
    Claim,
    Adjudication,
    Payment,
    PreAuthRequest,
    PreAuthDecision,
    Appeal,
    (*@\added{ProviderStatement,}@*)
    (*@\added{PlanAccumulator,}@*)
    (*@\added{AppealEvent,}@*)
    (*@\added{PharmacyClaim,}@*)
    (*@\added{MedicationCoverageCheck,}@*)
    (*@\added{PaymentPlan,}@*)
}
\end{lstlisting}

The \texttt{From<\&FactPayload> for FactType} implementation is updated
accordingly. The compiler enforces exhaustiveness, so every new
\texttt{FactPayload} variant requires a corresponding
\texttt{FactType} match arm.

\section{Migration and ingestion considerations}

The extensions in this document are additive with respect to the graph
model. Existing \texttt{Coverage}, \texttt{Claim}, \texttt{Adjudication},
\texttt{Payment}, \texttt{PreAuthDecision}, and \texttt{Appeal} facts
must be migrated to the extended structures.

Migration follows three guarantees: preserve the original representation;
do not invent source-authored specificity; and make every extraction,
reconstruction, or inference traceable to a source citation or prior Fact.

\begin{itemize}

\item \texttt{Claim.billed\_codes} and \texttt{Claim.billed\_amount}
migrate to a single-line \texttt{ClaimLine} where the source data does
not carry line detail. Ingestion pipelines that receive claim-level data
without lines produce a synthetic \texttt{ClaimLine} with the aggregate
amount and codes; the \texttt{line\_number} is \texttt{1} and the
enclosing Fact uses \texttt{AssertionMethod::Inferred}. Its evidence cites the
source record or prior Fact, and migration metadata records the reconstruction
rule version.

\item Existing \texttt{Coverage} facts without \texttt{effective\_from}
are migrated using \texttt{Fact.occurred\_at} as
\texttt{effective\_from}. \texttt{effective\_through} is set to
\texttt{None} for the most recent coverage and to the next coverage's
\texttt{effective\_from} for prior coverages, provided they are
non-overlapping.

\item Existing \texttt{PreAuthDecision.reason} strings are preserved in
the new \texttt{denial\_narrative} field. The
\texttt{denial\_codes} vector is left empty. A subsequent enrichment
pass may map free-text reasons to CARC/RARC codes, but it creates a separate
inferred coding Fact with the original denial Fact as evidence. It does not
supersede payer-authored content merely because enrichment became available.

\item Existing \texttt{Appeal} facts are preserved as-is. The
\texttt{appeal\_id} field is populated from any source-system
identifier available. If no such identifier exists, the field remains absent;
an internal deterministic key may support implementation but must not be
presented as a source identifier. \texttt{appeal\_level} defaults to
\texttt{Internal1}. \texttt{AppealEvent} facts are not
retroactively constructed; the appeal lifecycle is tracked
prospectively from the migration date forward.

\item Existing \texttt{Payment} facts are preserved as-is. The
\texttt{payment\_method} field defaults to \texttt{Other} where the
source data does not distinguish, and can be enriched during ingestion
from provider statement records if those are also ingested.

\end{itemize}

Ingestion pipelines from payer sources (CARIN Blue Button, Patient
Access APIs where available) produce the extended structures directly.
Ingestion from provider EHR exports (Epic EHI, generic FHIR) continues
to produce medical facts; billing detail from EHR exports is typically
limited to claim summaries and is ingested as reconstructed
single-line claims per above.

For every path, the original PDF, transaction, API response, or EHI export is
registered as a \texttt{SourceArtifact}; normalized Facts cite exact passages,
records, or rows. Source-provided claim and statement identifiers remain
\texttt{ExternalRef} values and are not confused with those citations.

\section{What this does not add}
\label{sec:what-not-added}

The extensions in this document capture billing and insurance events at
patient-facing granularity. They do not add:

\begin{itemize}

\item \textbf{Full plan design.} A \texttt{PlanDesign} object
describing formulary contents, service-level coverage rules, network
tiers, and benefit-year rules is not included. The minimal
\texttt{PlanStructure} type on \texttt{Coverage}
(Section~\ref{sec:plan-structure}) captures the plan's structural
category (HMO, PPO, HDHP, etc.) and cost-share dimensions (has
deductible, has copay, has coinsurance), but not the full plan
document. Individual accumulator state
(Section~\ref{sec:accumulator}) is captured, but the plan's benefit
schedule is not. If a benefits-intelligence use case emerges, full
plan design belongs in a separate contextual model outside
\texttt{FactPayload}.

\item \textbf{Vendor and benefit program utilization.} Enrollment in and
interaction with employer-sponsored benefit vendors (fertility, mental
health, MSK, diabetes management) is not modeled. These are relevant
to a benefits-intelligence product but not to a patient-facing health
record.

\item \textbf{Employer or plan-sponsor identity on coverage.}
\texttt{Coverage} identifies the payer and plan but not the plan
sponsor. Adding sponsor identity would enable aggregation across a
workforce but would make the FEN record employer-legible in ways that
conflict with the patient-owned positioning.

\item \textbf{Negotiated rate and market pricing context.} Cost
transparency comparisons (``what did this employee pay for a knee MRI
versus the plan's negotiated rate at another facility'') require
reference pricing data alongside claims. This is out of scope for
patient-facing billing defense and belongs in a separate pricing model
if pursued.

\end{itemize}

\section{Ontological placement summary}

\begin{center}
\small
\renewcommand{\arraystretch}{1.15}
\begin{tabular}{>{\raggedright\arraybackslash}p{5.4cm} >{\raggedright\arraybackslash}p{7.6cm}}
\textbf{Concern} & \textbf{Placement} \\
\hline
Original EOB, statement, export, or transaction & Immutable \texttt{SourceArtifact} outside FEN \\
Claim, Adjudication, Payment, Appeal, or billing event & FEN Fact payload \\
Payer, provider, pharmacy, or PBM identity & Contextual \texttt{Institution} registry \\
Claim-to-payment or statement-to-claim connection & Authored \texttt{FactRelation} \\
Evidence supporting a billing Fact & \texttt{Fact.evidence} \\
Source-system identity or correlation key & \texttt{Fact.external\_refs} \\
Current appeal state or reconciliation view & Projection or operational materialization \\
\end{tabular}
\end{center}

These omissions are deliberate. The extensions in this document
strengthen Phoros as a patient-owned health passport with meaningful
billing awareness. They do not turn Phoros into a benefits-intelligence
platform.

\end{document}
